\section{The logarithmic-time limit with finite preparation}
\label{sec:logtime}
We keep the state, group, equal independent matter species, exact energy,
fixed proper Gaussian width and the two threshold protocols of
Section~\ref{sec:crosslimit}. The limit is still strictly $G_N\to0$.
The result below is a distribution and packet limit in this model, not
a finite-gravity scrambling theorem. Throughout this section the repeated
times are $(t_M,0,t_M,0)$, with
\begin{equation}
 t_M=\log(2M)+T,\qquad T\in\mathbb R.
 \label{eq:logtime-scaling}
\end{equation}
The order of limits in the early and late plateaux below is first
$M\to\infty$ at fixed $T$, and only then $T\to\mp\infty$.

\subsection{Phase orientation and the surviving distribution}
For representation-label offsets write
$\chi=q_1-q_2+q_3-q_4$, $\delta=i\chi$, and use $p,q,e,t,v,U_\pm$
from Section~\ref{sec:assembly}; the auxiliary scalar $t$ there is not
the physical time $t_M$. Define
\begin{equation}
 V_+=2^{-\delta}U_+,\qquad V_-=2^{\delta}U_-.
 \label{eq:logtime-V}
\end{equation}
Both have value $2\pi$ on $\chi=0$. The leading expression (11.9) becomes
\begin{equation}
 \Lambda S^0_\Lambda=\frac{1}{\pi^3}
 \left[(2\Lambda)^{i\chi}\Gamma(i\chi)V_+
       +(2\Lambda)^{-i\chi}\Gamma(-i\chi)V_-\right].
 \label{eq:logtime-leading}
\end{equation}
Each summand on the right is assigned its symmetric principal value in
$\chi$; their sum is the original removable, ordinary function. This is
an algebraic decomposition, not an arbitrary choice of contact term.

In our ordering $(a,b,c,d)$ the repeated-time phase is
$e^{it_M\chi_m}$, where $\chi_m=m_a-m_b+m_c-m_d$; see
\eqref{eq:cross-time}. The minus branch, rather than the plus branch,
therefore survives positive logarithmic time. Set
\begin{equation}
 \mathcal L_T(q_1,q_2,q_3,q_4)
 =\frac{2}{\pi}\delta_D(\chi)
 +\frac{1}{\pi^3}\operatorname{PV}_{\chi}
      \left[e^{iT\chi}\Gamma(-i\chi)V_-\right].
 \label{eq:logtime-local-distribution}
\end{equation}
Here $\operatorname{PV}_\chi$ means symmetric exclusion of $|\chi|<h$
followed by $h\downarrow0$, at fixed three transverse linear coordinates.
The transverse choice does not alter this linear-hyperplane distribution.

\begin{theorem}[Local logarithmic-time limit]
On compact real-offset tests, for fixed $\nu>0$,
\begin{equation}
 \Lambda e^{i[\log(2\Lambda)+T]\chi}
 S(\Lambda+q_1,\ldots,\Lambda+q_4,\nu,\nu)
 \longrightarrow\mathcal L_T.
 \label{eq:logtime-raw-limit}
\end{equation}
Convergence is uniform for $T$ in compact real intervals.
\label{thm:logtime-local}
\end{theorem}
\begin{proof}
The remainder in Theorem~11.1, after multiplication by $\Lambda$ and a
unit-modulus phase, is $O(\Lambda^{-1+\epsilon})$ on each compact set.
Thus it remains to handle \eqref{eq:logtime-leading}. Locally,
\[
 \Gamma(i\chi)V_+=-\frac{2\pi i}{\chi}+R_+,
 \qquad
 \Gamma(-i\chi)V_-=\frac{2\pi i}{\chi}+R_-,
\]
with smooth $R_\pm$. Multiplication of the plus branch by the time phase
gives frequency $2\log(2\Lambda)+T$. The Fourier identity
\[
 e^{ik\chi}\operatorname{PV}\frac1\chi
 \longrightarrow i\pi\operatorname{sgn}(k)\delta_D(\chi),
 \qquad |k|\to\infty,
\]
and integration by parts for the smooth remainder give the positive
$2\delta_D(\chi)/\pi$ term. The minus branch has only the finite phase
$e^{iT\chi}$. These arguments are uniform on compact $T$ intervals:
the large frequency remains positive and the finite-phase test seminorms
are bounded there. Smooth compact transverse integrations cause no change.
\end{proof}

\subsection{Restoring external factors and exact energy}
Let $\xi=(x,y,z,0)$, $\chi=x-y+z$, and define the smooth factor
\begin{equation}
 H_\nu(x,y,z)=\frac{e^{-\pi y}}{1024\pi^5}
 \prod_{j=1}^{4}\prod_{\varepsilon=\pm1}
 \Gamma\left(\frac14+\frac i2
       [\xi_j-\xi_{j+1}+\varepsilon\nu]\right),
 \quad \xi_5=\xi_1.
 \label{eq:logtime-H}
\end{equation}
It is the $F_\Lambda/\Lambda$ limit of (11.14), with its complex phase
retained. In particular
$4H_\nu(x,x+z,z)/\pi=v_\nu(x)v_\nu(z)$.
Let $\mathcal L_T(x,y,z,0)$ denote the pullback of
\eqref{eq:logtime-local-distribution}. Translation of all four labels
does not change it. Then in the exact-energy variables of (9.7),
\begin{equation}
 e^{it_M(x-y+z)}\widetilde{\mathcal C}_M(r;x,y,z)
 \longrightarrow H_\nu(x,y,z)\mathcal L_T(x,y,z,0)
 \quad\hbox{in }\cD'(\mathbb R^4).
 \label{eq:logtime-energy-limit}
\end{equation}
To verify that the changing time introduces no missing constant, change
the four energy variables to their exact representation labels as in
Theorem~11.2. The smooth test and external-factor errors are $O(M^{-1})$
on a compact set. The difference between the two linear phase arguments
is at worst $O(M^{-1})$, so their exponentials differ by
$O(M^{-1}\log M)$, uniformly for compact $T$. The pointwise compact
bound $\Lambda S=O(1+\log\Lambda)$ bounds the resulting pairing error
by $O(M^{-1}\log^2M)$. The three internal Jacobians cancel in the change
of variables; the external inverse Jacobian tends to one. For energy
differences alone the sharper error is $O(M^{-2})$, but it is not needed.
These estimates also justify replacing $\Lambda$ by $M$ in the logarithm.

\subsection{Hard thresholds and arbitrary fixed packets}
Put $Q=x^2+(y-x)^2+(y-z)^2+z^2$, and let $K_0$ be (9.9). The two
limiting threshold weights after integrating the external energy are
\begin{equation}
 W_F(x,y,z)=1,\qquad
 W_Q(x,y,z)=K_0\bigl(\max(0,-x,-y,-z)\bigr).
 \label{eq:logtime-threshold-weights}
\end{equation}
Define
\begin{equation}
 \mathcal X_P(T)=
 \int_{\mathbb R^3}^{\rm reg}
  W_P(x,y,z)e^{-\sigma^2 Q/2}H_\nu(x,y,z)
           \mathcal L_T(x,y,z,0)\,\dd x\dd y\dd z,
 \qquad P=F,Q.
 \label{eq:logtime-packet}
\end{equation}
The superscript specifies the following definition: use smooth compact
external densities and path cutoffs, smooth every intermediate threshold,
pair the resulting smooth test with the displayed distribution, and then
remove these approximations with the bounds below. It does not silently
multiply a general distribution by a discontinuous function. For a general
$L^1$ density this definition is retained; no ordinary absolutely convergent
three-dimensional integral or unsmoothed pointwise principal value is asserted.

\begin{theorem}[Finite-packet limit and ordered plateaux]
For every fixed probability density $\pi_0\in L^1(\mathbb R_+)$ and either
protocol, the regularized expression \eqref{eq:logtime-packet} exists,
is independent of the allowed approximations, and
\begin{equation}
 X^P_M(t_M,0,t_M,0)\longrightarrow\mathcal X_P(T)
 \label{eq:logtime-packet-convergence}
\end{equation}
uniformly on compact $T$ intervals. It obeys
\begin{equation}
 \lim_{T\to-\infty}\mathcal X_P(T)=X^P_\infty,
 \qquad \lim_{T\to+\infty}\mathcal X_P(T)=0,
 \label{eq:logtime-plateaux}
\end{equation}
where $X^P_\infty$ is the repeated-time fixed-time limit in Section~9.
The limiting commutator norm squared is
\begin{equation}
 \mathcal D_P(T)=2\bigl(N^P_\infty-\operatorname{Re}\mathcal X_P(T)\bigr).
 \label{eq:logtime-commutator}
\end{equation}
Consequently $\mathcal D_F(-\infty)=0$,
$\mathcal D_Q(-\infty)=\mathcal D_{\rm edge}$, and
$\mathcal D_P(+\infty)=2N^P_\infty$.
\label{thm:logtime-packet}
\end{theorem}
\begin{proof}
All truncation estimates of Proposition~9.1 hold for every real time,
so also for $t_M$. Each threshold layer of width $\eta$ excludes a
tree norm squared at most $2\eta A\sup B$, giving a pairing error
$O(\sqrt\eta)$. Smooth path cutoffs have the same positive tail control.
Changing the external density in $L^1$ costs at most $A^2$ times its
$L^1$ distance. These bounds are independent of $M$ and $T$.
For each fixed smooth approximation, \eqref{eq:logtime-energy-limit}
gives convergence. The approximation estimates pass to its limit and
make the family of limiting pairings uniformly Cauchy for all $T$.
This proves existence, approximation independence and compact-$T$
uniform convergence. It also derives \eqref{eq:logtime-threshold-weights}
from $r,r+x,r+y,r+z\ge0$, rather than assuming the fixed-time delta
constraint away from its support.

For a fixed compact smooth test, use the same Fourier identity on the
minus branch of \eqref{eq:logtime-local-distribution}. Its singular part
contributes $-2\operatorname{sgn}(T)\delta_D(\chi)/\pi$, while its
smooth remainder vanishes. Thus $\mathcal L_T\to4\delta_D(\chi)/\pi$
as $T\to-\infty$ and $\mathcal L_T\to0$ as $T\to+\infty$.
The approximation bounds just proved are uniform in $T$, allowing their
removal after either limit. On $y=x+z$, the Gaussian, external factor
and threshold weight recover exactly the fixed-time expression (9.9).
The ordered norms are time independent and converge by the one-step
positive-kernel limit and its Gaussian envelope. Equation~\eqref{eq:commutator-norm}, or
direct expansion of the two-state norm, gives
\eqref{eq:logtime-commutator}.
\end{proof}
Neither monotonicity nor an early-time exponential rate follows. In
particular $\mathcal D_Q(T)-\mathcal D_{\rm edge}$ need not be positive
from this argument. A hard moving boundary changes the analytic test
class needed for contour shifts and cannot be ignored when extracting
an exponent. The fixed-protocol limit is independent of $\pi_0$;
the moving-protocol limit generally is not.

\subsection{The constant time origin in the source}
There is a useful direct coefficient check against KL (B.50), using
the same raw ordering for this comparison. Let
$a_\pm=q_1-q_2\pm\nu$ and define $R(z)=\Gamma(1-z)\sin(\pi z/2)$.
The duplication and reflection identities give $c(z)=2^zR(z)$, so
\begin{equation}
 V_+=2\pi\cosh(\pi\chi/2)
 \frac{R(p+v)R(q+v)}{R(p-t)R(q-t)}.
 \label{eq:logtime-origin-identity}
\end{equation}
Here $p+v=1/2+ia_-$, $q+v=1/2+ia_+$,
$p-t=1/2+i(a_--\chi)$ and $q-t=1/2+i(a_+-\chi)$.
Reflecting the minus-$\nu$ Gamma factors in this expression gives
exactly the coefficient multiplying $\Gamma(i\chi)$ in KL (B.50),
after division by $\pi^3$. This is an algebraic identity off the pole,
extended by meromorphy. For the source's negative phase convention,
the exact assembled limit therefore has the surviving factor
$e^{-iT\chi}\Gamma(i\chi)V_+/\pi^3$ at
$\tau=\log(2\Lambda)+T$.

KL (6.32) uses that time origin, but (B.45) uses
$\log(2\Lambda/e)$ with the same printed coefficient in (B.50).
At one and the same physical time, write
$T_{\rm app}=\tau-\log(2\Lambda/e)=T+1$.
Our surviving factor is then $e^{-i(T_{\rm app}-1)\chi}$, whereas the
printed (B.51) uses $e^{-iT_{\rm app}\chi}$. The difference is a finite
$e^{i\chi}$ phase. Source footnote 73 acknowledges an additive shift;
the calculation here fixes our convention from the exact Gamma ratios.
Exact-energy corrections vanish and cannot supply this constant. We
record this as a source-convention/asymptotic discrepancy, not an
author-confirmed erratum. It does not affect the fixed-time delta or
the two plateaux, but must be kept in any quantitative curve comparison.
